\documentclass[12pt,twoside]{article}
\usepackage{fancyhdr,makeidx,times,multicol,geometry}
\usepackage{keyval,ifthen,mparhack,manyfoot,poemscol}
\newcounter{scounter}
\newcommand\step{\stepcounter{scounter}\Roman{scounter}}

\title{Postal del Cielo}
\author{J.J. Gómez Cadenas}
\begin{document}
\maketitle

\poemtitle{Postcard from Heaven}
\begin{poem}
\begin{stanza}
But then, again,\\
you return to the city with your children.
\end{stanza}
\begin{stanza}
Your daughter is twenty-two, still holds your hand,\\
talks non-stop about her latest love-boy.\\
He is coming to meet her.\\
They will walk the same streets,\\
and kiss in the same corners\\
that once smiled at you.\\
They are inventing, one more time,\\
the same old tale,\\
yet so new.\\
The same old lies,\\
yet so true.
\end{stanza}
\begin{stanza}
Your son is eighteen,\\
wants to be a scientist like Mom and Dad.\\
When he was small,\\
he wanted to work in the same place\\
as both of you. Now\\
he is moving to Paris, away from home,\\
to fulfil his dreams.
\end{stanza}
\begin{stanza}
How strange, to feel so happy,\\
and so hopeless. So proud of him,\\
as a king watching his prince being knighted, yet\\
with a black hole eating your heart.
\end{stanza}
\begin{stanza}
Took so long to learn,\\
love was not all those will-o'-the-wisps,\\
vanishing in the night,\\
like the light of fireflies.
\end{stanza}
\begin{stanza}
Love was something else,\\
made of clay and ashes,\\
made of blood and tears,\\
made of guilt and forgiveness,\\
made of years.
\end{stanza}
\begin{stanza}
Took so long to learn,\\
Paris was not Paradise,\\
but it can still be\\
a postcard from Heaven.
\end{stanza}
\end{poem}

\poemtitle{Postal del Cielo}
\begin{poem}
\begin{stanza}
Pero entonces, de nuevo,\\
vuelves a la ciudad con tus hijos.
\end{stanza}
\begin{stanza}
Tu niña tiene ya veintidós años, aún te toma de la mano,\\
habla sin cesar de su nuevo enamorado.\\
Pronto se encontrarán en el mismo París,\\
que te conoció a su edad, andarán por las mismas calles,\\
y se besarán en las mismas esquinas,\\
que un día te sonrieron.\\
Inventando otra vez,\\
esa vieja historia, tan repetida,\\
y siempre nueva,\\
susurrando las viejas mentiras, tan verdaderas.
\end{stanza}
\begin{stanza}
Tu muchacho, con dieciocho,\\
quiere ser científico como sus padres.\\
De pequeño, soñaba trabajar con vosotros,\\
ahora sueña con volar y su vuelo,\\
le trae, precisamente, a París,\\
tan cerca y tan lejos a la vez.
\end{stanza}
\begin{stanza}
Qué extraño, sentirse a la vez,\\
tan feliz y tan desesperanzado.\\
Tan orgulloso de ellos,\\
como un rey que contempla la gloria de sus príncipes,\\
y sentir, al mismo tiempo,\\
un agujero negro formarse en tu pecho.
\end{stanza}
\begin{stanza}
Tanto tiempo para comprender,\\
que el amor no eran esos fuegos fatuos,\\
esfumándose en la noche,\\
como luz de luciérnagas.
\end{stanza}
\begin{stanza}
No, el amor era otra cosa,\\
hecho de barro y cenizas,\\
de sangre y lágrimas,\\
de culpa y perdón,\\
hecho de lentos años.
\end{stanza}
\begin{stanza}
Tanto tiempo para comprender,\\
que París no era el Paraíso,\\
pero aún puede ser,\\
una postal del cielo.
\end{stanza}
\end{poem}

\end{document}
